% comparison_EN.tex — WAC confronted with KAPE and Eric Zimmerman's tools, English edition.
% The French edition is comparison_FR.tex, in this folder: both are maintained together,
% section for section.
%
% Build (XeLaTeX, two passes for the cover page and the table of contents):
%     latexmk -xelatex comparison_EN.tex
% The cover page is ../cover_comparison_EN.png, derived from ../cover_EN.png by
% ../retitle-cover.py. The layout is ../wac-style.tex, shared with every WAC document.
%
% Every figure comes from backlog/B-015 and the runs of vmtest/ez-compare.sh and
% vmtest/kape-compare.sh, which reproduce them.

\documentclass[11pt,a4paper,oneside]{report}

% ---------------------------------------------------------------- language and labels
\usepackage[english]{babel}
\newcommand{\wacdoctitle}{Against KAPE and the EZ tools}
\newcommand{\wacpageof}{of}
\newcommand{\wacnotetitle}{Key point}
\newcommand{\wacwarningtitle}{Limit}
% Layout shared with every WAC document.
% wac-style.tex — layout shared by every WAC document (user manuals, testing
% guides), French and English editions alike: one copy, so that they cannot
% drift apart.
%
% A document sets, BEFORE % wac-style.tex — layout shared by every WAC document (user manuals, testing
% guides), French and English editions alike: one copy, so that they cannot
% drift apart.
%
% A document sets, BEFORE % wac-style.tex — layout shared by every WAC document (user manuals, testing
% guides), French and English editions alike: one copy, so that they cannot
% drift apart.
%
% A document sets, BEFORE \input{../wac-style.tex}:
%   \usepackage[<language>]{babel}
%   \newcommand{\wacdoctitle}{...}      running header: "WAC — <title>"
%   \newcommand{\wacpageof}{...}        footer: "Page N <of> M"
%   \newcommand{\wacnotetitle}{...}     default title of the "note" box
%   \newcommand{\wacwarningtitle}{...}  default title of the "attention" box
% and AFTER it its \hypersetup (PDF title, author, subject).

% ---------------------------------------------------------------- fonts, language
\usepackage{fontspec}
\setmainfont{Noto Sans}
\setmonofont{DejaVu Sans Mono}[Scale=0.84]

% ---------------------------------------------------------------- layout
\usepackage[top=26mm,bottom=26mm,left=24mm,right=22mm,headheight=14pt]{geometry}
\usepackage{graphicx}
\graphicspath{{../}}
\usepackage{xcolor}
\usepackage{tikz}
\usepackage{booktabs,tabularx,longtable,array}
\usepackage{enumitem}
\usepackage{fancyhdr}
\usepackage{titlesec}
\usepackage[most]{tcolorbox}
\usepackage{listings}
\usepackage{url}
\usepackage{lastpage}
\usepackage[hidelinks,unicode,pdfpagelabels]{hyperref}

% Colours sampled from the cover page
\definecolor{marine}{HTML}{0A1A33}
\definecolor{bleu}{HTML}{2F6DB5}
\definecolor{accent}{HTML}{1F66D1}
\definecolor{fond}{HTML}{EEF3FA}
\definecolor{alerte}{HTML}{B3261E}
\definecolor{fondalerte}{HTML}{FBEDEC}
\definecolor{gris}{HTML}{5B6573}

\setlength{\parindent}{0pt}
\setlength{\parskip}{0.55em}
\linespread{1.08}
\setlist{itemsep=0.25em,topsep=0.35em}
\urlstyle{tt}
\renewcommand{\arraystretch}{1.25}
% Left-aligned table columns: justified narrow columns opened wide gaps
% between words.
\newcolumntype{L}[1]{>{\raggedright\arraybackslash}p{#1}}
% Monospaced file and field names cannot be hyphenated: let the line's spaces
% stretch instead.
\setlength{\emergencystretch}{3em}

% Headings
\titleformat{\chapter}[display]
  {\color{marine}\bfseries\raggedright}
  {\color{bleu}\large\MakeUppercase{\chaptertitlename}~\thechapter}
  {0.3em}
  {\Huge}[{\color{accent}\titlerule[2pt]}]
\titlespacing*{\chapter}{0pt}{0pt}{1.6em}
\titleformat{\section}{\color{marine}\Large\bfseries}{\color{bleu}\thesection}{0.8em}{}
\titleformat{\subsection}{\color{marine}\large\bfseries}{\color{bleu}\thesubsection}{0.7em}{}
\titleformat{\subsubsection}{\color{bleu}\normalsize\bfseries}{}{0em}{}

% Headers and footers
\pagestyle{fancy}
\fancyhf{}
\fancyhead[L]{\small\color{gris}WAC — \wacdoctitle}
\fancyhead[R]{\small\color{gris}\leftmark}
% Footer: "Page N of M" (the word set by \wacpageof), readable. A bare, small, light-grey number went
% unnoticed.
\newcommand{\piedpage}{\small\color{marine}\textbf{Page \thepage}\ \wacpageof\ \pageref*{LastPage}}
\fancyfoot[C]{\piedpage}
\renewcommand{\footrulewidth}{0.4pt}
\renewcommand{\footrule}{\hbox to\headwidth{\color{fond}\leaders\hrule height \footrulewidth\hfill}}
\renewcommand{\headrulewidth}{0.4pt}
\renewcommand{\headrule}{\hbox to\headwidth{\color{fond}\leaders\hrule height \headrulewidth\hfill}}
\renewcommand{\chaptermark}[1]{\markboth{#1}{}}
\fancypagestyle{plain}{\fancyhf{}\fancyfoot[C]{\piedpage}\renewcommand{\headrulewidth}{0pt}}

% Boxes
\newtcolorbox{note}[1][\wacnotetitle]{enhanced,breakable,colback=fond,colframe=bleu,
  boxrule=0pt,leftrule=3pt,arc=0pt,left=8pt,right=8pt,top=6pt,bottom=6pt,
  title={\bfseries #1},coltitle=bleu,colbacktitle=fond,attach title to upper={\par\smallskip}}
\newtcolorbox{attention}[1][\wacwarningtitle]{enhanced,breakable,colback=fondalerte,colframe=alerte,
  boxrule=0pt,leftrule=3pt,arc=0pt,left=8pt,right=8pt,top=6pt,bottom=6pt,
  title={\bfseries #1},coltitle=alerte,colbacktitle=fondalerte,attach title to upper={\par\smallskip}}

% Code
\lstset{basicstyle=\ttfamily\small,columns=fullflexible,keepspaces=true,
  backgroundcolor=\color{fond},frame=l,framerule=3pt,rulecolor=\color{bleu},
  xleftmargin=9pt,framexleftmargin=6pt,aboveskip=0.8em,belowskip=0.8em,
  breaklines=true,breakatwhitespace=false,showstringspaces=false,
  literate={é}{{\'e}}1 {è}{{\`e}}1 {à}{{\`a}}1 {ê}{{\^e}}1 {ô}{{\^o}}1 {ç}{{\c c}}1
           {É}{{\'E}}1 {—}{{---}}1 {’}{{'}}1 {…}{{...}}1 {«}{{\guillemotleft}}1 {»}{{\guillemotright}}1}

% Breakable monospace. A path or a field name in \texttt is one word that no
% hyphenation pattern cuts: in a table cell it ran 190 pt past the column
% (ProgramData\Microsoft\NetFramework\BreadcrumbStore). As the \seqsplit of
% the EventsManager theme's tables, a line may break inside it, without a
% hyphen that would read as part of the name: after \ / . _ - : and between a
% lowercase letter or digit and a capital (AmcacheSha1|Covers|Bytes). Only
% inside \texttt — the classes are set once, the breaks switched on there.
\newXeTeXintercharclass\wacBreakAfter
\newXeTeXintercharclass\wacUpper
\newXeTeXintercharclass\wacLower
\XeTeXcharclass`\\=\wacBreakAfter
\XeTeXcharclass`\/=\wacBreakAfter
\XeTeXcharclass`\.=\wacBreakAfter
\XeTeXcharclass`\_=\wacBreakAfter
\XeTeXcharclass`\-=\wacBreakAfter
\XeTeXcharclass`\:=\wacBreakAfter
\count255=`A \loop\XeTeXcharclass\count255=\wacUpper\ifnum\count255<`Z\advance\count255 by 1\repeat
\count255=`a \loop\XeTeXcharclass\count255=\wacLower\ifnum\count255<`z\advance\count255 by 1\repeat
\count255=`0 \loop\XeTeXcharclass\count255=\wacLower\ifnum\count255<`9\advance\count255 by 1\repeat
\XeTeXinterchartoks\wacBreakAfter\wacUpper={\allowbreak}
\XeTeXinterchartoks\wacBreakAfter\wacLower={\allowbreak}
\XeTeXinterchartoks\wacBreakAfter 0={\allowbreak}
\XeTeXinterchartoks\wacLower\wacUpper={\allowbreak}
% \textbackslash is no character token: it carries its own break.
\newcommand{\wacbackslash}{\char92\allowbreak}
\let\wacplaintexttt\texttt
\renewcommand{\texttt}[1]{\wacplaintexttt{\XeTeXinterchartokenstate=1\let\textbackslash\wacbackslash#1}}

% Shortcuts
\newcommand{\wac}{\textsc{wac}}
\newcommand{\opt}[1]{\texttt{-{}-#1}}
\newcommand{\champ}[1]{\texttt{#1}}
\newcommand{\fichier}[1]{\texttt{#1}}

% ============================================================ COVER PAGE
% The cover images are 2:3 (1024 x 1536), taller than an A4 page once widened
% to it: centred, they lost 9 mm at the top — the CEC logo — and 9 mm at the
% bottom. Raised 4 mm above the top of the page instead: the logo, which starts
% 6 mm below the image's edge, stays whole, and the bottom row of text keeps
% 7 mm to the page's edge; only a decorative corner below it is cut.
% Usage: \coverpage{cover_FR.png}
\newcommand{\coverpage}[1]{%
  \begin{titlepage}
  \thispagestyle{empty}
  \begin{tikzpicture}[remember picture,overlay]
    \node[inner sep=0pt,anchor=north,yshift=4mm] at (current page.north)
      {\includegraphics[width=\paperwidth]{#1}};
  \end{tikzpicture}
  \end{titlepage}}
:
%   \usepackage[<language>]{babel}
%   \newcommand{\wacdoctitle}{...}      running header: "WAC — <title>"
%   \newcommand{\wacpageof}{...}        footer: "Page N <of> M"
%   \newcommand{\wacnotetitle}{...}     default title of the "note" box
%   \newcommand{\wacwarningtitle}{...}  default title of the "attention" box
% and AFTER it its \hypersetup (PDF title, author, subject).

% ---------------------------------------------------------------- fonts, language
\usepackage{fontspec}
\setmainfont{Noto Sans}
\setmonofont{DejaVu Sans Mono}[Scale=0.84]

% ---------------------------------------------------------------- layout
\usepackage[top=26mm,bottom=26mm,left=24mm,right=22mm,headheight=14pt]{geometry}
\usepackage{graphicx}
\graphicspath{{../}}
\usepackage{xcolor}
\usepackage{tikz}
\usepackage{booktabs,tabularx,longtable,array}
\usepackage{enumitem}
\usepackage{fancyhdr}
\usepackage{titlesec}
\usepackage[most]{tcolorbox}
\usepackage{listings}
\usepackage{url}
\usepackage{lastpage}
\usepackage[hidelinks,unicode,pdfpagelabels]{hyperref}

% Colours sampled from the cover page
\definecolor{marine}{HTML}{0A1A33}
\definecolor{bleu}{HTML}{2F6DB5}
\definecolor{accent}{HTML}{1F66D1}
\definecolor{fond}{HTML}{EEF3FA}
\definecolor{alerte}{HTML}{B3261E}
\definecolor{fondalerte}{HTML}{FBEDEC}
\definecolor{gris}{HTML}{5B6573}

\setlength{\parindent}{0pt}
\setlength{\parskip}{0.55em}
\linespread{1.08}
\setlist{itemsep=0.25em,topsep=0.35em}
\urlstyle{tt}
\renewcommand{\arraystretch}{1.25}
% Left-aligned table columns: justified narrow columns opened wide gaps
% between words.
\newcolumntype{L}[1]{>{\raggedright\arraybackslash}p{#1}}
% Monospaced file and field names cannot be hyphenated: let the line's spaces
% stretch instead.
\setlength{\emergencystretch}{3em}

% Headings
\titleformat{\chapter}[display]
  {\color{marine}\bfseries\raggedright}
  {\color{bleu}\large\MakeUppercase{\chaptertitlename}~\thechapter}
  {0.3em}
  {\Huge}[{\color{accent}\titlerule[2pt]}]
\titlespacing*{\chapter}{0pt}{0pt}{1.6em}
\titleformat{\section}{\color{marine}\Large\bfseries}{\color{bleu}\thesection}{0.8em}{}
\titleformat{\subsection}{\color{marine}\large\bfseries}{\color{bleu}\thesubsection}{0.7em}{}
\titleformat{\subsubsection}{\color{bleu}\normalsize\bfseries}{}{0em}{}

% Headers and footers
\pagestyle{fancy}
\fancyhf{}
\fancyhead[L]{\small\color{gris}WAC — \wacdoctitle}
\fancyhead[R]{\small\color{gris}\leftmark}
% Footer: "Page N of M" (the word set by \wacpageof), readable. A bare, small, light-grey number went
% unnoticed.
\newcommand{\piedpage}{\small\color{marine}\textbf{Page \thepage}\ \wacpageof\ \pageref*{LastPage}}
\fancyfoot[C]{\piedpage}
\renewcommand{\footrulewidth}{0.4pt}
\renewcommand{\footrule}{\hbox to\headwidth{\color{fond}\leaders\hrule height \footrulewidth\hfill}}
\renewcommand{\headrulewidth}{0.4pt}
\renewcommand{\headrule}{\hbox to\headwidth{\color{fond}\leaders\hrule height \headrulewidth\hfill}}
\renewcommand{\chaptermark}[1]{\markboth{#1}{}}
\fancypagestyle{plain}{\fancyhf{}\fancyfoot[C]{\piedpage}\renewcommand{\headrulewidth}{0pt}}

% Boxes
\newtcolorbox{note}[1][\wacnotetitle]{enhanced,breakable,colback=fond,colframe=bleu,
  boxrule=0pt,leftrule=3pt,arc=0pt,left=8pt,right=8pt,top=6pt,bottom=6pt,
  title={\bfseries #1},coltitle=bleu,colbacktitle=fond,attach title to upper={\par\smallskip}}
\newtcolorbox{attention}[1][\wacwarningtitle]{enhanced,breakable,colback=fondalerte,colframe=alerte,
  boxrule=0pt,leftrule=3pt,arc=0pt,left=8pt,right=8pt,top=6pt,bottom=6pt,
  title={\bfseries #1},coltitle=alerte,colbacktitle=fondalerte,attach title to upper={\par\smallskip}}

% Code
\lstset{basicstyle=\ttfamily\small,columns=fullflexible,keepspaces=true,
  backgroundcolor=\color{fond},frame=l,framerule=3pt,rulecolor=\color{bleu},
  xleftmargin=9pt,framexleftmargin=6pt,aboveskip=0.8em,belowskip=0.8em,
  breaklines=true,breakatwhitespace=false,showstringspaces=false,
  literate={é}{{\'e}}1 {è}{{\`e}}1 {à}{{\`a}}1 {ê}{{\^e}}1 {ô}{{\^o}}1 {ç}{{\c c}}1
           {É}{{\'E}}1 {—}{{---}}1 {’}{{'}}1 {…}{{...}}1 {«}{{\guillemotleft}}1 {»}{{\guillemotright}}1}

% Breakable monospace. A path or a field name in \texttt is one word that no
% hyphenation pattern cuts: in a table cell it ran 190 pt past the column
% (ProgramData\Microsoft\NetFramework\BreadcrumbStore). As the \seqsplit of
% the EventsManager theme's tables, a line may break inside it, without a
% hyphen that would read as part of the name: after \ / . _ - : and between a
% lowercase letter or digit and a capital (AmcacheSha1|Covers|Bytes). Only
% inside \texttt — the classes are set once, the breaks switched on there.
\newXeTeXintercharclass\wacBreakAfter
\newXeTeXintercharclass\wacUpper
\newXeTeXintercharclass\wacLower
\XeTeXcharclass`\\=\wacBreakAfter
\XeTeXcharclass`\/=\wacBreakAfter
\XeTeXcharclass`\.=\wacBreakAfter
\XeTeXcharclass`\_=\wacBreakAfter
\XeTeXcharclass`\-=\wacBreakAfter
\XeTeXcharclass`\:=\wacBreakAfter
\count255=`A \loop\XeTeXcharclass\count255=\wacUpper\ifnum\count255<`Z\advance\count255 by 1\repeat
\count255=`a \loop\XeTeXcharclass\count255=\wacLower\ifnum\count255<`z\advance\count255 by 1\repeat
\count255=`0 \loop\XeTeXcharclass\count255=\wacLower\ifnum\count255<`9\advance\count255 by 1\repeat
\XeTeXinterchartoks\wacBreakAfter\wacUpper={\allowbreak}
\XeTeXinterchartoks\wacBreakAfter\wacLower={\allowbreak}
\XeTeXinterchartoks\wacBreakAfter 0={\allowbreak}
\XeTeXinterchartoks\wacLower\wacUpper={\allowbreak}
% \textbackslash is no character token: it carries its own break.
\newcommand{\wacbackslash}{\char92\allowbreak}
\let\wacplaintexttt\texttt
\renewcommand{\texttt}[1]{\wacplaintexttt{\XeTeXinterchartokenstate=1\let\textbackslash\wacbackslash#1}}

% Shortcuts
\newcommand{\wac}{\textsc{wac}}
\newcommand{\opt}[1]{\texttt{-{}-#1}}
\newcommand{\champ}[1]{\texttt{#1}}
\newcommand{\fichier}[1]{\texttt{#1}}

% ============================================================ COVER PAGE
% The cover images are 2:3 (1024 x 1536), taller than an A4 page once widened
% to it: centred, they lost 9 mm at the top — the CEC logo — and 9 mm at the
% bottom. Raised 4 mm above the top of the page instead: the logo, which starts
% 6 mm below the image's edge, stays whole, and the bottom row of text keeps
% 7 mm to the page's edge; only a decorative corner below it is cut.
% Usage: \coverpage{cover_FR.png}
\newcommand{\coverpage}[1]{%
  \begin{titlepage}
  \thispagestyle{empty}
  \begin{tikzpicture}[remember picture,overlay]
    \node[inner sep=0pt,anchor=north,yshift=4mm] at (current page.north)
      {\includegraphics[width=\paperwidth]{#1}};
  \end{tikzpicture}
  \end{titlepage}}
:
%   \usepackage[<language>]{babel}
%   \newcommand{\wacdoctitle}{...}      running header: "WAC — <title>"
%   \newcommand{\wacpageof}{...}        footer: "Page N <of> M"
%   \newcommand{\wacnotetitle}{...}     default title of the "note" box
%   \newcommand{\wacwarningtitle}{...}  default title of the "attention" box
% and AFTER it its \hypersetup (PDF title, author, subject).

% ---------------------------------------------------------------- fonts, language
\usepackage{fontspec}
\setmainfont{Noto Sans}
\setmonofont{DejaVu Sans Mono}[Scale=0.84]

% ---------------------------------------------------------------- layout
\usepackage[top=26mm,bottom=26mm,left=24mm,right=22mm,headheight=14pt]{geometry}
\usepackage{graphicx}
\graphicspath{{../}}
\usepackage{xcolor}
\usepackage{tikz}
\usepackage{booktabs,tabularx,longtable,array}
\usepackage{enumitem}
\usepackage{fancyhdr}
\usepackage{titlesec}
\usepackage[most]{tcolorbox}
\usepackage{listings}
\usepackage{url}
\usepackage{lastpage}
\usepackage[hidelinks,unicode,pdfpagelabels]{hyperref}

% Colours sampled from the cover page
\definecolor{marine}{HTML}{0A1A33}
\definecolor{bleu}{HTML}{2F6DB5}
\definecolor{accent}{HTML}{1F66D1}
\definecolor{fond}{HTML}{EEF3FA}
\definecolor{alerte}{HTML}{B3261E}
\definecolor{fondalerte}{HTML}{FBEDEC}
\definecolor{gris}{HTML}{5B6573}

\setlength{\parindent}{0pt}
\setlength{\parskip}{0.55em}
\linespread{1.08}
\setlist{itemsep=0.25em,topsep=0.35em}
\urlstyle{tt}
\renewcommand{\arraystretch}{1.25}
% Left-aligned table columns: justified narrow columns opened wide gaps
% between words.
\newcolumntype{L}[1]{>{\raggedright\arraybackslash}p{#1}}
% Monospaced file and field names cannot be hyphenated: let the line's spaces
% stretch instead.
\setlength{\emergencystretch}{3em}

% Headings
\titleformat{\chapter}[display]
  {\color{marine}\bfseries\raggedright}
  {\color{bleu}\large\MakeUppercase{\chaptertitlename}~\thechapter}
  {0.3em}
  {\Huge}[{\color{accent}\titlerule[2pt]}]
\titlespacing*{\chapter}{0pt}{0pt}{1.6em}
\titleformat{\section}{\color{marine}\Large\bfseries}{\color{bleu}\thesection}{0.8em}{}
\titleformat{\subsection}{\color{marine}\large\bfseries}{\color{bleu}\thesubsection}{0.7em}{}
\titleformat{\subsubsection}{\color{bleu}\normalsize\bfseries}{}{0em}{}

% Headers and footers
\pagestyle{fancy}
\fancyhf{}
\fancyhead[L]{\small\color{gris}WAC — \wacdoctitle}
\fancyhead[R]{\small\color{gris}\leftmark}
% Footer: "Page N of M" (the word set by \wacpageof), readable. A bare, small, light-grey number went
% unnoticed.
\newcommand{\piedpage}{\small\color{marine}\textbf{Page \thepage}\ \wacpageof\ \pageref*{LastPage}}
\fancyfoot[C]{\piedpage}
\renewcommand{\footrulewidth}{0.4pt}
\renewcommand{\footrule}{\hbox to\headwidth{\color{fond}\leaders\hrule height \footrulewidth\hfill}}
\renewcommand{\headrulewidth}{0.4pt}
\renewcommand{\headrule}{\hbox to\headwidth{\color{fond}\leaders\hrule height \headrulewidth\hfill}}
\renewcommand{\chaptermark}[1]{\markboth{#1}{}}
\fancypagestyle{plain}{\fancyhf{}\fancyfoot[C]{\piedpage}\renewcommand{\headrulewidth}{0pt}}

% Boxes
\newtcolorbox{note}[1][\wacnotetitle]{enhanced,breakable,colback=fond,colframe=bleu,
  boxrule=0pt,leftrule=3pt,arc=0pt,left=8pt,right=8pt,top=6pt,bottom=6pt,
  title={\bfseries #1},coltitle=bleu,colbacktitle=fond,attach title to upper={\par\smallskip}}
\newtcolorbox{attention}[1][\wacwarningtitle]{enhanced,breakable,colback=fondalerte,colframe=alerte,
  boxrule=0pt,leftrule=3pt,arc=0pt,left=8pt,right=8pt,top=6pt,bottom=6pt,
  title={\bfseries #1},coltitle=alerte,colbacktitle=fondalerte,attach title to upper={\par\smallskip}}

% Code
\lstset{basicstyle=\ttfamily\small,columns=fullflexible,keepspaces=true,
  backgroundcolor=\color{fond},frame=l,framerule=3pt,rulecolor=\color{bleu},
  xleftmargin=9pt,framexleftmargin=6pt,aboveskip=0.8em,belowskip=0.8em,
  breaklines=true,breakatwhitespace=false,showstringspaces=false,
  literate={é}{{\'e}}1 {è}{{\`e}}1 {à}{{\`a}}1 {ê}{{\^e}}1 {ô}{{\^o}}1 {ç}{{\c c}}1
           {É}{{\'E}}1 {—}{{---}}1 {’}{{'}}1 {…}{{...}}1 {«}{{\guillemotleft}}1 {»}{{\guillemotright}}1}

% Breakable monospace. A path or a field name in \texttt is one word that no
% hyphenation pattern cuts: in a table cell it ran 190 pt past the column
% (ProgramData\Microsoft\NetFramework\BreadcrumbStore). As the \seqsplit of
% the EventsManager theme's tables, a line may break inside it, without a
% hyphen that would read as part of the name: after \ / . _ - : and between a
% lowercase letter or digit and a capital (AmcacheSha1|Covers|Bytes). Only
% inside \texttt — the classes are set once, the breaks switched on there.
\newXeTeXintercharclass\wacBreakAfter
\newXeTeXintercharclass\wacUpper
\newXeTeXintercharclass\wacLower
\XeTeXcharclass`\\=\wacBreakAfter
\XeTeXcharclass`\/=\wacBreakAfter
\XeTeXcharclass`\.=\wacBreakAfter
\XeTeXcharclass`\_=\wacBreakAfter
\XeTeXcharclass`\-=\wacBreakAfter
\XeTeXcharclass`\:=\wacBreakAfter
\count255=`A \loop\XeTeXcharclass\count255=\wacUpper\ifnum\count255<`Z\advance\count255 by 1\repeat
\count255=`a \loop\XeTeXcharclass\count255=\wacLower\ifnum\count255<`z\advance\count255 by 1\repeat
\count255=`0 \loop\XeTeXcharclass\count255=\wacLower\ifnum\count255<`9\advance\count255 by 1\repeat
\XeTeXinterchartoks\wacBreakAfter\wacUpper={\allowbreak}
\XeTeXinterchartoks\wacBreakAfter\wacLower={\allowbreak}
\XeTeXinterchartoks\wacBreakAfter 0={\allowbreak}
\XeTeXinterchartoks\wacLower\wacUpper={\allowbreak}
% \textbackslash is no character token: it carries its own break.
\newcommand{\wacbackslash}{\char92\allowbreak}
\let\wacplaintexttt\texttt
\renewcommand{\texttt}[1]{\wacplaintexttt{\XeTeXinterchartokenstate=1\let\textbackslash\wacbackslash#1}}

% Shortcuts
\newcommand{\wac}{\textsc{wac}}
\newcommand{\opt}[1]{\texttt{-{}-#1}}
\newcommand{\champ}[1]{\texttt{#1}}
\newcommand{\fichier}[1]{\texttt{#1}}

% ============================================================ COVER PAGE
% The cover images are 2:3 (1024 x 1536), taller than an A4 page once widened
% to it: centred, they lost 9 mm at the top — the CEC logo — and 9 mm at the
% bottom. Raised 4 mm above the top of the page instead: the logo, which starts
% 6 mm below the image's edge, stays whole, and the bottom row of text keeps
% 7 mm to the page's edge; only a decorative corner below it is cut.
% Usage: \coverpage{cover_FR.png}
\newcommand{\coverpage}[1]{%
  \begin{titlepage}
  \thispagestyle{empty}
  \begin{tikzpicture}[remember picture,overlay]
    \node[inner sep=0pt,anchor=north,yshift=4mm] at (current page.north)
      {\includegraphics[width=\paperwidth]{#1}};
  \end{tikzpicture}
  \end{titlepage}}


\hypersetup{pdftitle={WAC — Against KAPE and the EZ tools},
  pdfauthor={Centre d'Excellence Cyberdéfense Aérospatiale (CEC)},
  pdfsubject={WAC confronted with KAPE and Eric Zimmerman's tools: time, accuracy, formats}}

\begin{document}

% ============================================================ COVER PAGE
\pagenumbering{roman}
\coverpage{cover_comparison_EN.png}

% ============================================================ INFORMATION PAGE
\thispagestyle{empty}
\vspace*{\fill}
{\color{marine}\Large\bfseries WAC — Windows Artefacts Collector}\par
{\color{bleu}Against KAPE and Eric Zimmerman's tools}\par
\medskip
\begin{tabular}{@{}ll@{}}
Software version & 1.4.0 \\
Document edition & October 2026 \\
Publisher & Centre d'Excellence Cyberdéfense Aérospatiale (CEC) \\
          & École de l'air et de l'espace, air base 701, Salon-de-Provence \\
\end{tabular}

\bigskip
{\small\color{gris}This document confronts \wac{} with the community's
reference tools: KAPE for collection, Eric Zimmerman's command-line tools for
decoding. On the same exhibits of the same machine, it compares running times,
what each one takes, the accuracy of the values --- every difference settled by
an independent reference ---, the output formats, and the traces each one leaves
on the examined machine. It says what the confrontation taught
\wac{}, and what it still holds against it.}
\vspace*{2cm}
\clearpage

\pagenumbering{arabic}
\tableofcontents

% ============================================================ 1
\chapter{The method}

\section{What is compared}

Two questions, two comparisons.

\begin{description}[style=nextline,leftmargin=1em]
  \item[Collection: \wac{} against KAPE] Both tools collect the same machine
        one after the other; what each copies is compared, the time it takes,
        and the SHA-1 of the files both copy.
  \item[Traces left: \wac{}, KAPE and the EZ tools] Each one is started from the
        same state of the same machine, under Sysmon; what it writes there is
        recorded, what Windows writes there because of it, and what a
        monitoring tool sees of it (chapter~\ref{chap:traces}).
  \item[Decoding: \wac{} against the EZ tools] On the \emph{same exhibits} ---
        the exhibit store of a \wac{} collection ---, \wac{} converts an artefact
        and the matching Eric Zimmerman tool decodes it; both outputs are compared
        record by record and field by field.
\end{description}

\section{Who is right when they differ}

A difference between two tools does not say which one is wrong. Every difference
was therefore settled by a third source, independent of both:

\begin{itemize}[nosep]
  \item Windows itself: the XML it renders for an event (\texttt{Get-WinEvent}),
        a file's date, the NTFS object identifier (\texttt{fsutil objectid}), the
        live registry, a file's hash (\texttt{Get-FileHash});
  \item another implementation: \texttt{dissect.esedb} (Fox-IT) for ESE
        databases;
  \item the format itself: sequence numbers and Marvin32 digests of the hive
        logs, read byte by byte;
  \item two artefacts of the same machine cross-checked: Prefetch for execution,
        a row's timestamp for its other dates.
\end{itemize}

Each difference is then classed: \emph{defect of \wac{}} (fixed), \emph{error of
the EZ tool} (checked), or \emph{convention} of either --- two forms of one value.

\section{The test bench}

\begin{tabular}{@{}ll@{}}
Machine & Windows 11 Pro (build 26200), virtual machine \\
EZ tools & version 2026.5.0 (.NET 9.0.20) \\
KAPE & version 1.3.0.2, target \texttt{KapeTriage} \\
\wac{} & version 1.4.0 with the fixes this confrontation brought \\
Traces measure & clone of the machine, no network, frozen; Sysmon 15.22; \\
               & tools on a USB key (K:), one run per tool \\
\end{tabular}

Everything can be replayed: \fichier{vmtest/ez-compare.sh <artefact>} times
\wac{} and the EZ tool on the same exhibits then compares their outputs;
\fichier{vmtest/kape-compare.sh} does the same for collection;
\fichier{vmtest/footprint/run-footprint.sh} replays the measure of the traces
and \fichier{footprint-diff.py} analyses it. The EZ tools ran
with their full date precision (\opt{mp}, or \opt{dt} with seven decimals where
\opt{mp} is not enough), so as not to blame them for a loss of precision a
setting avoids.

% ============================================================ 2
\chapter{Summary}

\begin{longtable}{@{}L{3.2cm}L{2.4cm}L{4.5cm}L{5.0cm}@{}}
\toprule
\textbf{Artefact} & \textbf{Volume} & \textbf{Result} & \textbf{Errors of the EZ tool, checked} \\
\midrule\endhead
\$MFT --- MFTECmd & 362,081 records & paths, states, 8 dates identical & sizes of 1,490 WOF files and 3 open logs: Windows gives \wac{}'s \\
EVTX --- EvtxECmd & 190,613 events & system fields identical; values explained & 6 wrong values (Windows); 335 Perflib events unreadable; cleared logs ignored \\
Prefetch --- PECmd & 379 files & all identical & hash without its leading zero (19) \\
Shortcuts --- LECmd & 3 files & identical to 100\,ns & --- \\
Jump lists --- JLECmd & 13 shortcuts & identical to 100\,ns & --- \\
Shellbags --- SBECmd & 27 folders & identical & wrong NodeSlot at first level (7); HasExplored always false \\
Amcache --- AmcacheParser & 1,258 entries & 124 columns identical & 2,340 invented values; partial SHA-1 not flagged \\
ShimCache --- AppCompat\-Cache\-Parser & 232 entries & identical & misleading execution flag \\
Recycle bin --- RBCmd & 19 files & identical & --- \\
SRUM --- SrumECmd & 60,804 rows & identical & 17 wrong energy rows (dissect) \\
Registry --- RECmd & 9 artefacts & identical & 4 hives rolled back to an older state \\
Collection --- KAPE & 1,173 common files & SHA-1 identical & --- \\
\bottomrule
\end{longtable}

\begin{note}
On everything both sides read, \wac{} gives the same values as the reference
tools, or more accurate ones, each time established by an independent source.
The confrontation also revealed 24 defects of \wac{}, all fixed
(chapter~\ref{chap:defects}), and coverage gaps, opened for follow-up
(section~\ref{sec:coverage}).

Started from a key, \wac{} writes nothing on the examined machine: no file, no
registry value, no event --- only the traces Windows leaves of any program
started remain (Amcache, ShimCache). KAPE and the EZ tools, .NET programs,
write to the user's registry and, for KAPE's modules, files on the system
volume (chapter~\ref{chap:traces}).
\end{note}

% ============================================================ 3
\chapter{Running times}

\section{Decoding}

Same machine, same exhibits; each tool handles its artefact only.

\begin{longtable}{@{}L{3.6cm}L{2.2cm}L{2.2cm}L{7.0cm}@{}}
\toprule
\textbf{Artefact} & \textbf{\wac{}} & \textbf{EZ tool} & \textbf{What \wac{} does besides} \\
\midrule\endhead
\$MFT (380\,MB) & 7.3\,s & 8.1\,s & two volumes (C: and G:) against one \\
EVTX (305 logs) & 16.8\,s & 22.0\,s & plain-text message rebuilt for every event \\
Registry (Kroll batch) & 2.0\,s & 7.6 to 13.6\,s & every hive replayed, every registry artefact \\
SRUM (21\,MB) & 3.2\,s & 2.5\,s & 2 more tables, identities resolved \\
Prefetch (380) & 3.6\,s & 0.7\,s & dated local time, paths with drive letter \\
Shortcuts, jump lists & 1.8\,s & 0.3\,s & \\
Recycle bin & 1.9\,s & 0.5\,s & \\
Amcache, ShimCache, shellbags & 2.0\,s & 0.5 to 1.4\,s & one pass for all three (registry key) \\
\bottomrule
\end{longtable}

On the large artefacts, \wac{} is faster. On the small ones, the EZ tools are:
every \wac{} conversion pays a fixed cost of 1.5 to 3.4\,s, measured with a
trivial artefact. That cost is the verified working copy of the hives, the
replay of their logs, and the reading of the suspect's time zone and profiles.
It is what lets \wac{} date every local time with the offset in force at that
date, and give the drive letter; the EZ tools give neither.

\section{Collection}

\begin{longtable}{@{}L{5.4cm}L{2.2cm}L{2.4cm}L{5.0cm}@{}}
\toprule
\textbf{Tool, setting} & \textbf{Time} & \textbf{Exhibits} & \textbf{Volume} \\
\midrule\endhead
KAPE, \texttt{KapeTriage} & 58\,s & 1,047 & 0.97\,GB \\
\wac{}, kit, no binaries & 47\,s & 1,898 & 2.54\,GB \\
\wac{}, kit, no binaries, before B-033 & 212\,s & 2,622 & 1.69\,GB \\
\wac{}, full kit (\texttt{binary: unverified}), before B-033 & 454\,s & 61,294 & 3.6\,GB \\
\bottomrule
\end{longtable}

The first two rows are measured in the same run, one after the other. On a
comparable scope, \wac{} now collects faster than KAPE, bringing back more than
twice as much data. It was five times slower at first (item B-033):
\begin{itemize}
  \item each raw reading opened its volume again and walked its paths from the
        root; the volumes are now opened once, with the cache of their
        directories;
  \item the collection took the files of the 1,426 declared event providers
        (739\,MB, 61\,s); it now takes only those of the providers the
        collected logs cite (324 on the VM, 178\,MB, 15\,s), which is enough to
        rebuild every event's text.
\end{itemize}
The full kit also checks the Authenticode signature of every executable on the
machine and copies those whose signature does not hold, which KAPE does not do.

\section{Collection from one same state}

Measures of the traces comparison (chapter~\ref{chap:traces}): every tool
starts from the same frozen state of the machine, started from the key; the
duration is that of the task that starts it, in brackets the one the tool
reports.

\begin{longtable}{@{}L{5.6cm}L{2.6cm}L{3.4cm}L{3.4cm}@{}}
\toprule
\textbf{Tool, setting} & \textbf{Duration} & \textbf{Taken} & \textbf{On the key} \\
\midrule\endhead
KAPE, \texttt{!SANS\_Triage} & 127\,s (115\,s) & 2,291 files & 2,286 files, 1.25\,GB \\
KAPE, \texttt{!SANS\_Triage} and \texttt{!EZParser} modules & 235\,s (223\,s) & 2,288 files, 14 tools & 2,345 files, 1.63\,GB \\
EZ tools on the live machine (6 tools) & 132\,s & \$MFT, logs, hives, Prefetch, Amcache, ShimCache & 36 CSV, 0.55\,GB \\
\wac{}, collection (default setting) & 592\,s (588\,s) & 11,830 exhibits, 67,882 recorded & 8,240 files, 4.9\,GB \\
\bottomrule
\end{longtable}

\wac{}'s collection is the longest because it reads \emph{every} executable of
the machine: the conversion will happen elsewhere, and it does not know yet
which ones the artefacts will cite. It computes their fingerprints, checks their
signature and copies only those whose signature does not hold; hence 67,882
exhibits recorded for 11,830 copied. The conversion, on the analysis
workstation, reads nothing of the machine (chapter~\ref{chap:traces}). KAPE copies the files of its targets without
examining them.

% ============================================================ 4
\chapter{What each one takes}\label{chap:taken}

\section{The same exhibits, different files}

Same runs as the previous section: \wac{}'s collection, KAPE with
\texttt{!SANS\_Triage}, compared by the original path of every exhibit
(\wac{}'s manifest, KAPE's copy log).

\begin{longtable}{@{}L{6.4cm}L{8.6cm}@{}}
\toprule
\textbf{Exhibits} & \textbf{Number} \\
\midrule\endhead
Taken by both & 1,151, 1,011 of them with the same SHA-1 \\
Taken by KAPE alone & 1,140: the user's profile (930, browsers and
  applications), the NetworkService profile (91), \fichier{SleepStudy} (79),
  \fichier{LogFiles} (17), \fichier{WDI} (16) \\
Taken by \wac{} alone & 10,679: executables whose signature does not hold
  (5,300 of them copies the test bench left in \fichier{C:\textbackslash
  wactest}), signature catalogs (1,205), event logs and their providers'
  resource files, Defender (98) \\
\bottomrule
\end{longtable}

The 140 common exhibits with a different fingerprint are files Windows keeps
rewriting --- event logs (71), Prefetch (38), hives (6), SRUM (3): both runs
start from the same state, but Windows changed them since start-up,
differently each time. Their creation date is identical in 1,146 cases out of
1,151. KAPE's own coverage is the one backlog items B-016, B-031 and B-032
already follow.

\section{What is recorded of each exhibit}

\begin{longtable}{@{}L{3.6cm}L{6.0cm}L{5.4cm}@{}}
\toprule
 & \textbf{\wac{} (manifest)} & \textbf{KAPE (copy log)} \\
\midrule\endhead
Fingerprints & MD5, SHA-1, SHA-256 & SHA-1 \\
File dates & the 4 of \texttt{\$STANDARD\_INFORMATION} and the 4 of
  \texttt{\$FILE\_NAME}, to the 100\,ns, in UTC and in the suspect's local
  time & 3 of \texttt{\$STANDARD\_INFORMATION} (creation, modification,
  access), to the 100\,ns, in UTC, without the MFT record change \\
Location & MFT record number & path alone \\
Method & stated for every exhibit (raw NTFS reading, live observation) &
  ordinary copy; ``deferred'' raw reading for locked files (181 out of 2,292) \\
Failure & code and cause in the manifest, exhibit by exhibit & console and log \\
Integrity & sealed manifest (SHA-256), checked again at every conversion & --- \\
\bottomrule
\end{longtable}

\begin{note}
The \texttt{\$FILE\_NAME} dates are those an intruder does not rewrite with the
ordinary functions: a backdated file (``timestomping'') keeps its real creation
date in \texttt{\$FILE\_NAME}. KAPE's log does not have them; \wac{}'s manifest
gives them for every exhibit, with the MFT record that finds it again in
\fichier{mft.json}.
\end{note}

% ============================================================ 5
\chapter{Traces left on the examined machine}\label{chap:traces}

\section{The measure}

A clone of the test machine, without network, is frozen with Sysmon installed;
the tools are on a USB key, drive K:. Every run starts from that same state: the
system disk and the key are fresh layers over the frozen images. Once started
up, each tool is launched the same way --- a scheduled task of the user, with its
elevated rights, as an examiner from an elevated prompt.

\begin{description}[style=nextline,leftmargin=1em]
  \item[Files] The USN of C:'s change journal is read just before and just after
        the tool: every file created, changed, renamed or deleted between the two
        is in it, even one created then deleted during the run. The readings are
        written to the key, after the end bound.
  \item[Author] Sysmon gives the tool's process tree and what each process
        wrote, files and registry. A write of another process that names the
        tool (Amcache, Prefetch, BAM) is Windows reacting to its start.
  \item[Control] A run that starts nothing (\texttt{ver}) measures the launcher
        and Windows' own activity; what it shows too is not laid on the tool.
  \item[Deferred] The ShimCache is written at shutdown only: every run ends with
        a clean shutdown, a restart and its reading.
\end{description}

\section{What each one leaves}

\begin{longtable}{@{}L{2.5cm}L{3.0cm}L{4.0cm}L{5.5cm}@{}}
\toprule
\textbf{Tool} & \textbf{Files written on C:} & \textbf{Registry written} &
  \textbf{Windows' reactions} \\
\midrule\endhead
\wac{} (collection) & none & none & Amcache: one entry during the
  run; ShimCache: one entry at shutdown \\
KAPE, targets & none & 4 Internet Settings \texttt{ZoneMap} values in the
  user's hive; 169 open-or-create calls (certificate stores, policies) &
  Amcache and ShimCache for \fichier{kape.exe} \\
KAPE, targets and modules & 28: 14 in
  \fichier{ProgramData\textbackslash Microsoft\textbackslash NetFramework\textbackslash
  BreadcrumbStore}, 14 in the user's Temp
  (\fichier{system-commandline-sentinel-files}) & the same 4 values, rewritten
  by 14 of the 15 processes (136 writes); about 2,600 open-or-create calls &
  Amcache and ShimCache for KAPE and 14 tools \\
EZ tools, live & none & no value; 54 open-or-create calls (policies,
  \texttt{Tcpip}, \texttt{Internet Settings}) & Amcache and ShimCache for the
  6 tools \\
\bottomrule
\end{longtable}

No process of any tool made Windows write an event in its logs: those of the
measure window are the services' logons, Defender, group policy, common to the
control run or unrelated to the tool. All of them read the volumes raw, which
Sysmon records (\texttt{RawAccessRead}) and an EDR may report.

\begin{description}[style=nextline,leftmargin=1em]
  \item[Open-or-create] \texttt{RegCreateKeyEx} creates the key if it is missing
        and opens it otherwise; Sysmon does not tell the two apart. The .NET
        libraries open the certificate stores and the policy keys that way. A
        created key changes date; a key merely opened does not.
  \item[The \texttt{ZoneMap} values] Written by WinINet's start-up in every .NET
        process that loads it, they change the examined user's
        \fichier{NTUSER.DAT} hive, hence its log and the key's date.
  \item[The .NET files] The runtime leaves one ``breadcrumb'' per assembly
        loaded in \fichier{BreadcrumbStore}, and the System.CommandLine library a
        file per tool in the user's Temp. Started alone, the EZ tools left none;
        started by KAPE, they did.
  \item[Prefetch and BAM] None of the 22 programs started from the USB key got
        a Prefetch file or a BAM value, while \fichier{cmd.exe}, started from the
        system disk by the same task, has its BAM value. A tool copied onto the
        machine and started from there would get them.
  \item[The .NET prerequisite] KAPE and the EZ tools need the .NET~9 runtime,
        present on the clone. On a machine without it, installing it writes
        hundreds of files and keys; \wac{} is a self-contained executable.
\end{description}

\section{What the measure fixed in \wac{}}

\wac{} cleared the screen at start-up with \texttt{system("cls")}, which starts
\fichier{cmd.exe}: a process created on the examined machine and the Prefetch of
\fichier{CMD.EXE} updated, at every collection. The screen is now cleared
through the console API, \wac{} imports no function able to start a process any
more, and the build script refuses a version that would import one.

\begin{note}
Every trace of \wac{} is explained: it starts no process, writes no file on the
machine, no registry value, no event. What remains is what Windows leaves of any
program started --- an Amcache entry, a ShimCache entry at shutdown ---,
unavoidable and described in the harmlessness file, and the raw reading of the
volumes, visible to an EDR.
\end{note}

% ============================================================ 6
\chapter{Accuracy}

\section{In favour of \wac{}}

\begin{description}[style=nextline,leftmargin=1em]
  \item[Hive logs replayed in order] The \fichier{.LOG1} and \fichier{.LOG2} logs
        keep, behind the recent entries, leftovers of an earlier generation.
        RECmd applies them in file order, hence after the newer ones. It thus
        rolls SYSTEM, SECURITY, DEFAULT and a service profile's hive back to an
        older state: in SYSTEM, sequence 0x3618 instead of 0x3620, and four last
        BAM executions shown 13 minutes early. \wac{} sorts the entries by
        sequence and stops at the break in the chain. The logs themselves confirm
        it, Marvin32 digests included.
  \item[Events of cleared logs] A cleared log keeps its records: \wac{} reads
        them again and flags them \champ{EvtRecovered}, never presenting them as
        live; EvtxECmd does not see them.
  \item[Event values as Windows renders them] Hexadecimal \texttt{0x3e7}, GUIDs
        in braces, booleans \texttt{true}, binary \texttt{69060000}, dates with
        their \texttt{Z}, U+200E marks kept: the forms Event Viewer shows and one
        searches for. EvtxECmd changes eight of them, and is wrong on six values,
        such as a boolean \texttt{False} that Windows renders \texttt{true}.
  \item[Nothing invented] A value absent from the source is not emitted.
        AmcacheParser writes 2,340 default values absent from the hive
        (\texttt{IsPeFile: False} for a PE executable), SrumECmd 7,086 (0 or
        $-1$).
  \item[Partial Amcache SHA-1 flagged] For a file larger than 30\,MiB, Amcache
        hashes only the first 31,457,280 bytes (measured on \fichier{mrt.exe}).
        \wac{} says so (\champ{Sha1CoversBytes}); without it, comparing with the
        whole file's SHA-1 would wrongly conclude to a modification.
  \item[No misleading field] The ShimCache ``executed'' flag, which
        AppCompatCacheParser gives, says ``no'' for 105 of the 117 binaries that
        Prefetch proves were run. \wac{} does not collect it: that ``no'' would
        lead to conclude a program never ran.
  \item[Precision to 100\,ns] \wac{} gives every date at its source's
        precision. By default, PECmd, LECmd and SrumECmd round to the second, and
        LECmd ignores \opt{mp} in its CSV.
  \item[Dated local time] Every date is given in UTC and in the suspect's local
        time, with the offset in force at that date. The EZ tools give UTC only,
        without saying so by a \texttt{Z}.
  \item[Others' readings checked] SBECmd's ``FirstInteracted'' and
        ``LastInteracted'', RECmd's ``CreatedOn'' of an account, are readings of
        key write times. \wac{} publishes those times themselves: anyone can draw
        their own reading from them, without inheriting a ready-made one.
        SBECmd also gives a wrong \texttt{NodeSlot} (0) to first-level folders.
  \item[SrumECmd's ESE reading errors] 17 energy rows out of 887 carry the
        content of another row or an identifier absent from the table;
        dissect.esedb sides with \wac{}.
  \item[Collection] Each \wac{} exhibit carries MD5, SHA-1 and SHA-256, the
        eight NTFS dates of its original and its MFT entry, in a sealed manifest.
        KAPE gives one SHA-1 and three dates.
\end{description}

\section{In favour of the EZ tools}\label{sec:coverage}

\begin{description}[style=nextline,leftmargin=1em]
  \item[Registry coverage] RECmd's Kroll batch reads keys \wac{} ignores:
        installed programs, Defender exclusions, RecentDocs, MountPoints2, known
        networks, the NTFS last-access update setting, shutdown time, among
        others. Item B-031.
  \item[File coverage] Against \wac{}, KapeTriage now only copies the browsers
        and the Windows timeline, planned (B-016, B-010), and the
        \fichier{desktop.ini} files, left out. The shortcuts outside Recent, the
        Default profile's hive, \fichier{\$Secure:\$SDS}, \fichier{\$Boot},
        \fichier{\$UsnJrnl:\$Max}, \fichier{\$TxfLog}, Defender's WPP traces
        and the DPAPI master keys, which it alone took, are collected and
        converted since item B-032.
  \item[Speed on small artefacts and in collection] See the previous chapter.
  \item[Ready-to-read interpretations] The EZ tools split and name some fields:
        executable, date and description of a SRUM application, known folder
        names, a file's extension. \wac{} gives the raw value, as Windows stores
        it.
\end{description}

\begin{attention}
What \wac{} does not collect, it cannot weigh: KAPE's coverage remains wider on
browsers and several registry keys, until B-016, B-031 and B-032 are done.
\end{attention}

% ============================================================ 7
\chapter{Output formats}

\section{What each tool writes}

\begin{longtable}{@{}L{4.4cm}L{5.4cm}L{5.2cm}@{}}
\toprule
\textbf{Tool} & \textbf{Formats} & \textbf{Note} \\
\midrule\endhead
AmcacheParser & CSV & 8 files, one per table \\
AppCompatCacheParser & CSV & \\
RBCmd & CSV & \\
SrumECmd & CSV & 7 files \\
EvtxECmd & CSV, JSON, XML & \\
JLECmd & CSV, JSON, HTML & \\
LECmd & CSV, JSON, HTML, XML & \\
MFTECmd & CSV, JSON, bodyfile & \\
PECmd & CSV, JSON, HTML & \\
RECmd & CSV, JSON & plus one CSV per plugin \\
SBECmd & CSV, JSON & \\
KAPE & copy of the files, CSV copy log & \\
\wac{} & JSON for every artefact & one file per artefact, plus the sealed manifest \\
\bottomrule
\end{longtable}

\section{What sets them apart}

Four tools out of eleven write CSV only. In the others, the JSON transcribes the
CSV line: one text line per record (JSON Lines), flat.

\begin{description}[style=nextline,leftmargin=1em]
  \item[Types] In the EZ tools, numbers are strings: a run count is written
        \texttt{"RunCount":"458"}, a record number
        \texttt{"EventRecordId":"21572"}. \wac{} writes numbers, booleans,
        strings.
  \item[Structures] In the EZ tools, a list becomes a comma-joined string (the
        files a Prefetch loaded, its directories), and an event's data are JSON
        escaped \emph{inside} a string (\texttt{Payload}), to decode a second
        time. \wac{} nests: a Prefetch's volumes and their directories, a
        shortcut's items, an event's \champ{Name}/\champ{Value} pairs.
  \item[Dates] In the EZ tools, UTC is not flagged
        (\texttt{2026-09-30 23:04:46}) and precision depends on an option. \wac{}
        gives full ISO 8601, in UTC (\texttt{…Z}) and in local time with its
        offset (\texttt{…+02:00}).
  \item[Missing value] The EZ tools write a default value, or an empty cell.
        \wac{} omits the field; an artefact not collected is declared so
        (\texttt{NotRequested}, \texttt{NotCollected}), so that an empty list is
        not taken for an absence of traces.
  \item[Provenance] \wac{}'s manifest ties every value to the exhibit it comes
        from, with its hashes and original dates.
  \item[Console] The EZ tools detail their work on the console; some write the
        decoding itself there when no file is asked for. \wac{} writes only its
        progress there: everything it finds is in the files.
\end{description}

\begin{note}
For an analyst, the difference shows at the first processing: a \wac{} JSON
loads as it is into an analysis tool or into WAE, with nothing to split again;
an EZ output is handled like a spreadsheet, whose numbers must be converted back,
lists split again and event data decoded again.
\end{note}

% ============================================================ 8
\chapter{What the confrontation taught \wac{}}\label{chap:defects}

Each defect below produced a valid JSON that was wrong, or incomplete. All are
fixed, with a check that would have caught them.

\begin{longtable}{@{}L{3.0cm}L{12.0cm}@{}}
\toprule
\textbf{Artefact} & \textbf{Defect fixed} \\
\midrule\endhead
\$MFT & stale extension records merged with the base \\
EVTX & floats cut to 6 digits; binary as ``69 06 00 00''; hexadecimal integers and GUIDs in the wrong form; ProcessingErrorData ignored (21 events without data); events of cleared logs not read \\
Shortcuts & URI glued to the block after it; reading 2 bytes off without an ID list; ExtraData blocks not read (machine, tracking); target size not read; null DROIDs emitted \\
Jump lists & open count and MRU rank not read; DestList format version 1 not supported; application name corrupted \\
Shellbags & folder lost with everything below it; the folder's key time, rank and NodeSlot missing; ``Home'' folder without a name \\
Amcache & 2 tables out of 14 and 4 to 6 values read: no SHA-1, size, driver or device \\
ShimCache & tabs of packaged applications replaced by spaces \\
SRUM & FILETIME dates left as numbers; UTF-16 text read as ANSI \\
Registry & BAM limited to accounts with a profile (a deleted account lost its executions); UserAssist foreground time not read \\
Start-up & \texttt{system("cls")} started \fichier{cmd.exe} on the examined machine (traces measure) \\
Test harness & a crashing check silently stopped every following one \\
\bottomrule
\end{longtable}

% ============================================================ 9
\chapter{Limits of the comparison}

\begin{itemize}
  \item One machine, virtual, running Windows 11: some artefacts are few there
        (3 shortcuts in Recent, 2 USB keys). The formats are the same on a real
        machine, but not every rare case is represented.
  \item The two collections follow each other instead of being simultaneous: the
        files written between them (\$MFT, logs, Prefetch) rightly differ.
  \item The EZ tools ran with their maximum precision options; by default, they
        lose more.
  \item The traces are measured on a clone without network, one run per tool: a
        write of Windows the control run did not show is not laid on the tool
        unless Sysmon gives it to one of its processes; it stays listed as not
        attributed in \fichier{footprint-diff.py}'s results.
  \item Process creation auditing (Security 4688) is not active on the test
        machine: where it is, every tool leaves an event per process --- one for
        \wac{}, fifteen for KAPE and its modules.
  \item The comparison covers the common artefacts. What only one tool reads has
        no reference here: its accuracy rests on the checks of the validation
        report.
\end{itemize}

\end{document}
